\scp{\tsc{Nuclear Sula-Buru}}{09-01}
Within \tsc{Sula-Buru} there is a primary split between \ili{Ambelau}
and all other languages; \tsc{Nuclear Sula-Buru}.
\tsc{Nuclear Sula-Buru} is defined by *z/*y > **y,
*b > \it{f}, and **mb > \it{b}.

Examples of *z/*y > **y are shown in \trf{09-03}.
There is a secondary development of **y > {\0} in \ili{Lisela} and \ili{Buru}.
\ili{Ambelau} has *z > \it{l}
and *y > {\0} often with traces on the vowels.
Additional examples of *z/*y > **y > {\0} in \ili{Buru} include *layaR
‘sail’ > \it{laa} ‘sail’,
\it{laa-n} ‘dorsal fin’, *ma-tazəm > \it{em-tae} ‘sharp’,
and *tuzuq ‘point at,
point out’ > \it{tuu-k} ‘lift, raise; designate, appoint’.


\begin{table}\tcp{\tsc{Nuclear Sula-Buru} *z/*y > **y}{09-03}
%\begin{adjustbox}{width=\textwidth}
	\begin{threeparttable}\st{3.25pt}
	\begin{tabular}{lllllll}\lsptoprule
*gloss	&	big	&	crocodile	&	rain	&	chin, jaw	&	path	&	far\\
PMP	&	*Ra\tbr{y}a	&	*buqa\tbr{y}a\su{†}	&	*qu\tbr{z}an	&	*qa\tbr{z}ay	&	*\tbr{z}alan	&	*\tbr{z}auq\\ \midrule
\ili{Mangole}	&	\cg{(fanini)}	&	{\hr}fo\tbr{y}a	&	{\hr}u\tbr{y}a	&	\cg{(lamida lila)}	&	\cg{(sanafa)}	&	{\hr}\tbr{y}au\\
\ili{Sanana}	&	{\hr}a\tbr{y}a	&	{\hr}fua\tbr{y}a	&	{\hr}u\tbr{y}a	&	{\hr}n-a\tbr{y}o hiha	&	{\hr}\tbr{y}aa	&	{\hr}\tbr{y}au\\
\ili{Lisela}	&	\cg{(bagu-t)}	&	{\hr}ubaa	&	\cg{(dekat)}	&	{\hr}aa-n	&	\cg{(tuhu-n)}	&	\cg{(eb-rema-n)}\\
\ili{Buru}	&	{\hr}haa	&	\cg{(emhalat)}	&	\cg{(dekat)}	&	{\hr}aa-n	&	\cg{(toho-n)}	&	\cg{(b-rema-n)}\\
\dshmidrule{7}
\ili{Ambelau}	&	{\hr}e-hea	&	{\hr}buaa	&	{\hr}u\tbr{l}a, u\tbr{l}an	&	{\hr}a\tbr{l}a-mu	&	{\hr}\tbr{l}ah-lea	&	\cg{(em-rahae)}\\
		\lspbottomrule
	\end{tabular}
		\begin{tablenotes}\tnsize
			\item[†] {\ili{Mangole} \it{foya} is possibly a loan from \tsc{Taliabu}.
								\ili{Lisela} \it{ubaa} is a loan from \ili{Kayeli}.}
		\end{tablenotes}
	\end{threeparttable}
%\end{adjustbox}
\end{table}

Examples of *b > \it{f} in \tsc{Nuclear Sula-Buru} are given in \trf{09-04}.
\ili{Ambelau} has *b > \it{b, β, h}.
Additional examples include *bulan ‘moon’ (\trf{07-01}),
*bətaw ‘sister (m.s.)’ (\trf{07-03}),
and *baqəRu ‘new’ (\trf{07-11}).

\begin{table}\tcp{\tsc{Nuclear Sula-Buru} *b > **f}{09-04}
%\begin{adjustbox}{width=\textwidth}
	\st{3pt}
	\begin{tabular}{llllllll}\lsptoprule
*gloss	&	stone	&	pig	&	star	&	fruit	&	body hair	&	wake (tr.)	&	woman\\
PMP	&	*\tbr{b}atu	&	*\tbr{b}a\tbr{b}uy	&	*\tbr{b}ituqən	&	*\tbr{b}uaq	&	*\tbr{b}ulu	&	*\tbr{b}aŋun	&	*\tbr{b}\<in\>ahi\\ \midrule
\ili{Mangole}	&	{\hr}\tbr{f}atu	&	{\hr}\tbr{f}a\tbr{f}i	&	{\hr}\tbr{f}antui	&	{\hr}\tbr{f}ua	&	{\hr}\tbr{f}oo	&	\cg{(gega ol)}	&	{\hr}\tbr{f}ina\\
\ili{Sanana}	&	{\hr}\tbr{f}at	&	{\hr}\tbr{f}a\tbr{f}	&	{\hr}\tbr{f}atui	&	{\hr}\tbr{f}ua	&	{\hr}\tbr{f}oa	&		&	{\hr}\tbr{f}ina\\
\ili{Lisela}	&	\cg{(miat)}	&	{\hr}\tbr{ɸ}a\tbr{ɸ}u	&	\cg{(gae)}	&	{\hr}\tbr{ɸ}ua-n	&	{\hr}\tbr{ɸ}ulu-n	&	{\hr}\tbr{ɸ}aŋo	&	{\hr}\tbr{ɸ}ina\\
\ili{Buru}	&	{\hr}\tbr{f}atu	&	{\hr}\tbr{f}a\tbr{f}u	&	\cg{(tolot)}	&	{\hr}\tbr{f}ua-n	&	{\hr}\tbr{f}olo-n	&	{\hr}\tbr{f}aŋo	&	{\hr}ana\tbr{f}ina\\
\dshmidrule{8}
\ili{Ambelau}	&	{\hr}\tbr{b}aru	&	{\hr}\tbr{b}a\tbr{β}u	&	\cg{(emralai)}	&	{\hr}o\tbr{b}ua-ni	&	{\hr}\tbr{b}olu-e	&	{\hr}\tbr{p}anu	&	{\hr}e-\tbr{β}liɲa\\
		\lspbottomrule
	\end{tabular}
%\end{adjustbox}
\end{table}

The behaviour of PMP *uCu in \tsc{Nuclear Sula-Buru} languages warrants some discussion.
\ili{Buru} and \ili{Lisela} have a restriction in harmony of height with
back vowels separated by a consonant,
such that in a sample of over 7,000 unique lexemes,
there are no morpheme internal sequences of *[oCu],
or *[uCo], except in loans (\citealp[62]{Grimes1991c}, \citealp{GrimesGrimes2020}).
The different patterns are:

\begin{itemize}[noitemsep]
	\item	PMP *uCu > \it{uCu} in \ili{Lisela}
	\item	PMP *uCu > \it{oCo}, \it{uCu} in \ili{Buru} (with \it{oCo} being the dominant pattern)
	\item	PMP *uCu > \it{oCu} in \ili{Mangole} \citep[6f]{Bloyd2015}
	\item	PMP *uCu > \it{oCa} in \ili{Sanana} \citep[136f]{Bloyd2020}
\end{itemize}

The \ili{Ambelau} data (see \srf{09-02}) show PMP *uCu > \it{oCo,
oCu, uCu}, but are unclear due a scarcity of data,
and identifiable loans from \ili{Buru} and \ili{Kayeli}.
Examples of *uCu are given in \trf{09-05}.

\begin{table}\tcp{\tsc{Sula-Buru} *uCu(C){\#}}{09-05}
\begin{adjustbox}{width=\textwidth}
	\begin{threeparttable}\st{2pt}
	\begin{tabular}{lllllllll}\lsptoprule
*gloss	&	head	&	descend	&	burn	&	pound	&	fish trap	&	joint	&	louse	&	br\ili{east}\\
PMP	&	*q\tbr{u}l\tbr{u}	&	*t\tbr{u}R\tbr{u}n	&	*t\tbr{u}n\tbr{u}	&	*t\tbr{u}kt\tbr{u}k	&	*b\tbr{u}b\tbr{u}	&	*b\tbr{u}k\tbr{u}	&	*k\tbr{u}t\tbr{u}	&	*s\tbr{u}s\tbr{u}\\ \midrule
\ili{Mangole}	&	\cg{(ŋapu)}	&		&	{\hr}d\tbr{o}n\tbr{u}	&	{\hr}d\tbr{o}t\tbr{u}	&	{\hr}f\tbr{o}f\tbr{u}	&	{\hr}f\tbr{o}k\tbr{u}	&	{\hr}k\tbr{o}t\tbr{u}	&	{\hr}s\tbr{o}s\tbr{u}\\
\ili{Sanana}	&	\cg{(nap)}	&	\cg{(neu)}	&	{\hr}d\tbr{o}n\tbr{a}	&	{\hr}d\tbr{o}t\tbr{a}	&	{\hr}f\tbr{o}f\tbr{a}	&	{\hr}f\tbr{o}k\tbr{a}	&	{\hr}k\tbr{o}t\tbr{a}	&	{\hr}s\tbr{o}s\tbr{a}\\
\ili{Lisela}	&	{\hr}\tbr{u}l\tbr{u}-n	&	{\hr}t\tbr{u}h\tbr{u}	&	\cg{(sigi)}	&	{\hr}t\tbr{u}t\tbr{u}	&		&	{\hr}ɸ\tbr{u}k\tbr{u}-n	&	\cg{(yemin)}	&	\cg{(hono-n)}\su{†}\\
\ili{Buru}	&	{\hr}\tbr{o}l\tbr{o}-n	&	{\hr}t\tbr{o}h\tbr{o}	&	\cg{(pefa)}	&	{\hr}t\tbr{o}t\tbr{o}	&	{\hr}f\tbr{o}f\tbr{o}	&	{\hr}f\tbr{o}k\tbr{o}-n	&	{\hr}k\tbr{o}t\tbr{o}	&	{\hr}s\tbr{o}s\tbr{o}-n\\
\noalign{\vskip-1.5ex}
\multicolumn{9}{c}{\dashedline{125mm}} \\
\noalign{\vskip-0.5ex}
\ili{Ambelau}	&	{\hr}\tbr{o}l\tbr{o}-ni	&	{\hr}r\tbr{o}h\tbr{o}	&	\cg{(peβa)}	&	{\hr}t\tbr{o}r\tbr{o}	&	{\hr}b\tbr{u}β\tbr{u}	&		&	{\hr}ʔ\tbr{u}r\tbr{u}	&	{\hr}sis\tbr{u}-a\\
		\lspbottomrule
	\end{tabular}
		\begin{tablenotes}\tnsize
			\item[†] {One \ili{Lisela} word list has \it{siso-ni} ‘br\ili{east}’
								which shows irregular initial *u > \it{i}, like \ili{Ambelau}.}
		\end{tablenotes}
	\end{threeparttable}
\end{adjustbox}
\end{table}

There is a further development in \ili{Sanana}
and \ili{Mangole} of **l > {\0} /Vα{\gap}Vα (between like vowels).
But there is also an ordering of processes,
such that this \it{l}-deletion between like vowels had to happen
before final *u > \it{a} in \ili{Sanana}.
For example, PMP *bulu ‘hair, feather’ > \ili{Mangole} \it{foo},
\ili{Sanana} \it{foa}, \ili{Lisela} \it{ɸulu-n},
\ili{Buru} \it{folo-n}; PMP *puluq ‘group of ten’ > \ili{Mangole} \it{poo},
\ili{Sanana} \it{poa}, \ili{Buru} \it{polo}.
With great consistency, PMP *d > \it{r} in \ili{Buru}
and \ili{Lisela}, but *d > \it{l} in \ili{Mangole}
and *d > \it{h} in \ili{Sanana},
as shown in \trf{09-06}.
The \ili{Mangole} and \ili{Sanana} reflexes could be secondary developments from intermediate **r.
\ili{Ambelau} has *d > \it{l,
h} (see \srf{09-02} below).
The different reflexes of *d
and *z in \tsc{Nuclear Sula-Buru} means that \tsc{Sula-Buru}
does not participate in the merger of *d/*z which occurs in nearly
all \tsc{\ili{Ambon}-Seram} and \tsc{Seram-Tanimbar-Bomberai}
languages (\srf{07-03-03}).

\newpage
\begin{table}\tcp{\tsc{Nuclear Sula-Buru} *d}{09-06}
\begin{adjustbox}{width=\textwidth}
	\begin{threeparttable}\st{2pt}
	\begin{tabular}{lllllllll}\lsptoprule
*gloss	&	thorn\su{†}	&	prawn	&	near	&	upriver	&	stand	&	catch, seize	&	body filth	&	eye\\
PMP	&	*\tbr{d}uRi	&	*qu\tbr{d}aŋ	&	*a\tbr{d}ani, *\tbr{d}aŋi	&	*\tbr{d}aya	&	*kə\tbr{d}əŋ	&	*\tbr{d}akəp	&	*\tbr{d}aki	&	**\tbr{d}ama\\ \midrule
\ili{Mangole}	&	{\hr}\tbr{l}oi	&	{\hr}ua\su{‡}	&	{\hr}\tbr{l}ani	&	{\hr}\tbr{l}ai	&	{\hr}ge\tbr{l}i	&		&	{\hr}\tbr{l}aki	&	{\hr\hr}\tbr{l}ama\\
\ili{Sanana}	&	{\hr}\tbr{h}oi	&	{\hr}u\tbr{h}a	&	{\hr}\tbr{h}an	&	{\hr}\tbr{h}ai	&	{\hr}ge\tbr{h}i	&	{\hr}\tbr{h}ak kot	&		&	{\hr\hr}\tbr{h}ama\\
\ili{Lisela}	&	{\hr}\tbr{r}ohi-n	&	{\hr}u\tbr{r}an	&	{\hr}b-\tbr{r}ani-n	&	{\hr}\tbr{d}ae\su{‖}	&	{\hr}ke\tbr{r}e	&		&	{\hr}\tbr{r}aki	&	{\hr\hr}\tbr{r}ama-n\\
\ili{Buru}	&	{\hr}\tbr{r}ohi-n	&	{\hr}u\tbr{r}an\su{Δ}	&	{\hr}eb-\tbr{r}aŋi-n	&	{\hr}\tbr{d}ae\su{‖}	&	{\hr}ke\tbr{r}e-k	&	{\hr}ef-\tbr{r}ake	&	{\hr}\tbr{r}aki	&	{\hr\hr}\tbr{r}ama-n\\
\noalign{\vskip-1.5ex}
\multicolumn{9}{c}{\dashedline{141mm}} \\
\noalign{\vskip-0.5ex}
\ili{Ambelau}	&	\cg{(kukuni)}	&{\hr}u\tbr{l}an	&	\cg{(rapate)}	&		&	\cg{(βati)}	&		&		&	{\hr\hr}\tbr{l}ama-ni\\
		\lspbottomrule
	\end{tabular}
		\begin{tablenotes}\tnsize
			\item[†]	{\ili{Buru} and \ili{Lisela} \it{rohi-n} = ‘bone’.}
			\item[‡]	{\ili{Mangole} \it{ua} `prawn' has irregular *d > {\0}.}
			\item[Δ]	{\ili{Buru} \it{uran} is restricted to saltwater shrimp
								and lobster; \it{sehe} is used for freshwater shrimp
								and crayfish (taboo driven; compare \it{sehe} ‘move backwards,
								retreat’; see \citealp{GrimesMaryott1994}).}
			\item[‖]	{\ili{Buru} and \ili{Lisela} \it{dae} ‘upstream,
								towards the interior, towards an emic centre’ is a likely loan
								from an unknown source.
								The regular reflex would be unattested *[raa].}
		\end{tablenotes}
	\end{threeparttable}
\end{adjustbox}
\end{table}

There are three innovations in the lexicon shared by \tsc{Nuclear Sula-Buru}:

\begin{itemize}[noitemsep]
	\item	**muyu ‘no,
				not’; \ili{Mangole} \it{moyu,} \ili{Sanana} \it{moya},
				\ili{Buru}, \ili{Lisela} \it{moo} (\ili{Ambelau} \it{raenu})
	\item	**muya ‘jungle’; \ili{Mangole} \it{fiŋa muya} ‘fly (n.)’,
				\ili{Sanana} \it{fina muya} ‘fly (n.)’ (Compare \ili{Buru} \it{feŋa} ‘fly
				(n.)’, \it{feŋa mua} ‘horsefly, deerfly,
				large stinging fly.’), \ili{Buru},
				\ili{Lisela} \it{mua} ‘jungle, forest,
				bush’ (see \citet[633f]{GrimesGrimes2020} for full definition
				and collocations) 
	\item	**munda ‘cool breeze’; \ili{Mangole} \it{moru} ‘wind’,
				\ili{Sanana} \it{mora} ‘wind, air’,
				\ili{Lisela} \it{moda,} \ili{Buru} \it{moda} ‘cool,
				refreshing (v.i.); wind, air (n.)’ (see Grimes
				and \citet[625f]{GrimesGrimes2020} for full definition
				and collocations) 
\end{itemize}

\tsc{Nuclear Sula-Buru} divides into \tsc{Sula} (\ili{Mangole}
and \ili{Sanana}), and \tsc{Buru} (\ili{Buru} and \ili{Lisela}).
In addition to differences in the historical phonology,
the differences in the lexicon are significant with only 27--33{\%}
lexical similarity between five varieties of \tsc{Buru}
and two varieties of \tsc{Sula} \citep{Grimes2000a}.
